\clearpage
\section{WebUI 设计}
\label{sec:webui}
\subsection{设计目标与信息架构}
WebUI 是嘟嘟哒面向管理员的日常工作台。其设计围绕三个连续任务展开：掌握真实群聊与账号状态，为目标群配置智能体行为，检查一次回答的执行过程。界面同时承载 QQ 通信与智能体管理，因此始终保留当前账号、会话及连接状态，让操作对象与结果可以对应。

信息架构分为会话工作区和管理页面。会话工作区让管理员在阅读群聊时直接打开 Agent；管理页面集中处理联系人、请求通知、群服务、人工内测与模型接入。路由保留账号、会话类型和对象标识，刷新页面后仍能定位同一会话。

\begin{table}[H]\centering
\caption{WebUI 的信息架构与主要任务}
\begin{tabularx}{\textwidth}{@{}L{29mm}L{42mm}Y@{}}\toprule
入口 & 管理员任务 & 主要信息与操作\\\midrule
消息工作区 & 阅读和处理当前讨论 & 账号切换、会话列表、消息历史、搜索、输入与回复；可并排展开 Agent。\\
Agent 工作台 & 配置与检查群聊智能体 & 对话、运行、配置三个页签；展示本轮上下文、回答、实际模型和能力调用。\\
联系人与通知 & 定位人员、群与待处理事项 & 群与好友目录、群管理入口、请求及事件；操作绑定具体 QQ 账号。\\
群服务 & 选择群服务初值 & 服务档案、当前状态与操作结果；依赖管理会话及对应控制面后端。\\
人工内测 & 检查既有样本和回答质量 & 浏览已配置的脱敏样本、生成候选并登记人工反馈。\\
API Key 池与设置 & 维护模型接入及工作区 & 三档模型连接、Key 状态、应用到 Runtime；主题、缓存和连接相关设置。\\\bottomrule
\end{tabularx}
\end{table}

\subsection{工作区布局与视觉交互}
桌面工作区采用四列布局：账号与导航栏、会话列表、消息正文、Agent 面板。消息与 Agent 并排，减少管理员在讨论内容和配置结果之间切换的成本；Agent 可折叠，为长消息与资料留出空间。管理页面使用账号栏加主体内容的布局，顶部明确显示执行操作的账号。

\begin{figure}[H]\centering
\begin{tikzpicture}[x=1cm,y=1cm,font=\small]
\draw[draw=rulegray,fill=white,rounded corners=1pt] (0,0) rectangle (15.2,6.2);
\fill[panel] (0.01,5.3) rectangle (15.19,6.19);
\draw[rulegray] (0,5.3)--(15.2,5.3);
\node[anchor=west,font=\small\sffamily\bfseries,text=accent] at (0.25,5.76) {当前账号 / 会话};
\node[anchor=east,text=muted] at (14.95,5.76) {QQ 账号状态\quad Agent 状态};
\fill[panel] (0.01,0.01) rectangle (1.65,5.29);
\draw[rulegray] (1.65,0)--(1.65,5.3);
\draw[rulegray] (5,0)--(5,5.3);
\draw[rulegray] (10.15,0)--(10.15,5.3);
\node[align=center,text width=1.3cm] at (0.825,2.8) {账号\\[4mm]消息\\联系人\\通知\\[4mm]管理\\设置};
\node[anchor=north,align=left,text width=2.8cm] at (3.325,5.02) {\textbf{会话列表}\\[3mm]搜索 / 筛选\\[5mm]选中的群聊\\[3mm]其他群与私聊\\[3mm]最新消息与未读};
\node[anchor=north,align=left,text width=4.55cm] at (7.575,5.02) {\textbf{消息正文}\\[3mm]当前讨论与时间\\[3mm]发言者、引用与消息段\\[3mm]历史分页与搜索};
\draw[rulegray] (5,1.18)--(10.15,1.18);
\node[align=center,text=muted] at (7.575,0.58) {消息输入\quad 提及 / 回复 / 发送};
\node[anchor=north,align=left,text width=4.4cm] at (12.675,5.02) {\textbf{Agent 工作台}\\[2mm]对话\quad 运行\quad 配置\\[3mm]同群近期上下文\\[2mm]候选回答与资料来源\\[2mm]模型、能力与耗时};
\draw[rulegray] (10.15,1.18)--(15.2,1.18);
\node[align=center,text=muted] at (12.675,0.58) {问题输入\quad 运行预览};
\end{tikzpicture}

\caption{桌面工作区的结构示意。图中仅表达布局和信息职责，不使用真实账号或聊天截图。}
\label{fig:webui-layout}
\end{figure}

响应式布局围绕任务切换调整：中等宽度下，Agent 变为可收起的侧边面板；宽度不超过 860 px 时，会话列表、消息和 Agent 以单页切换呈现，底部保留主要导航，其余入口收进“更多”。页面保存选中的会话，避免从 Agent 返回聊天时丢失位置。

界面以低饱和青色标识主要操作，使用浅色背景、细分隔线和有限的状态色组织密集信息，并提供深色主题。在线、警示与失败状态同时给出文字，模型档位和结果类型保持明确标签。键盘焦点可见，搜索和编辑弹窗支持 Esc 关闭、Tab 循环及关闭后恢复焦点；这些细节支持连续管理操作。

\subsection{前后端分工与消息同步}
前端以 Vue 3 和 TypeScript 实现，\code{App.vue} 组合账号栏、会话列表、聊天窗和 Agent 面板；工作区状态由组合式逻辑管理，目录查询与页面数据通过服务适配器进入组件。类型与响应校验集中在数据边界，组件负责呈现、输入和交互反馈。

Node.js Gateway 向浏览器提供同源 HTTP 与 SSE。HTTP 处理工作区读取、历史分页、配置和明确触发的操作；SSE 推送消息及连接变化。网关经认证的 OneBot v11 反向 WebSocket 接收 NapCat 事件，并在服务端执行限定的协议动作。OneBot 连接凭据由网关保管。

\begin{figure}[H]\centering
\begin{tikzpicture}[node distance=8mm and 8mm]
\node[box,text width=13.5cm] (ui) {Vue 浏览器：会话工作区、Agent 工作台与管理页面};
\node[box,text width=13.5cm,below=of ui] (gateway) {Node.js Gateway\\工作区快照 / SSE 事件 / 历史读取 / 配置与预览请求};
\node[box,text width=6.1cm,below=11mm of gateway.south west,anchor=north west] (qq) {NapCat / OneBot v11\\账号、消息与 QQ 操作};
\node[box,text width=6.1cm,below=11mm of gateway.south east,anchor=north east] (runtime) {AstrBot / Dududa Runtime\\感知、模型、校园工具与回答};
\draw[<->,>=Stealth,draw=accent,thick] (ui)--node[right,labeltext]{同源 HTTP / SSE}(gateway);
\draw[<->,>=Stealth,draw=accent,thick] (gateway.south -| qq.north)--node[left,labeltext]{反向 WebSocket}(qq.north);
\draw[<->,>=Stealth,draw=accent,thick] (gateway.south -| runtime.north)--node[right,labeltext]{Runtime HTTP API}(runtime.north);
\end{tikzpicture}

\caption{WebUI 的数据与调用关系。QQ 同步和 Agent 预览共用网关入口，分别连接各自的后端。}
\label{fig:webui-architecture}
\end{figure}

首屏先读取已有工作区快照，联系人与历史在后台继续加载；慢联系人或慢 Agent 请求不会阻止会话框架显示。SSE 重连根据事件位置补发可重放事件，超出保存范围时刷新工作区。历史分页与账号、会话绑定，消息按规范化标识合并。

浏览器使用 Dexie/IndexedDB 缓存已取得的消息、会话、历史覆盖范围和草稿，主键包含账号与会话标识。搜索明确标注当前浏览器缓存范围；未取得的离线消息不被描述为已经同步。缓存清理作用于本地副本，QQ 数据仍以 NapCat 返回结果为依据。

\subsection{Agent 对话、运行与配置}
Agent 面板将使用过程拆成三个页签：\textbf{对话}接收管理员请求并呈现回答与资料；\textbf{运行}展示实际模型、推理与篇幅配置、使用的能力、耗时、上下文覆盖和错误原因；\textbf{配置}维护本群的行为偏好、插件开关与参与策略。一次请求的运行信息与回答保持关联，管理员可以从结果直接追查执行过程。

在当前正式 Runtime 预览路径中，请求依次经过以下步骤：
\begin{enumerate}
\item 浏览器提交当前 \code{accountId}、\code{conversationId} 和问题。
\item 网关校验输入与会话范围，从服务端读取同群近期历史；读取失败时直接说明原因。
\item 网关使用服务端凭据调用 AstrBot 的 Runtime 预览接口，由既有核心完成感知、路由、必要查询和回答组织。
\item Runtime 返回回答、模型、能力、耗时和上下文用量，网关校验后交给页面呈现。
\item 管理员在“运行”页查看此次实际执行情况，再决定是否调整配置或继续提问。
\end{enumerate}

预览模式可执行校园工具查询，但结果留在 Web 中，返回的 QQ 输出调用与长期记忆写入均为 0。QQ 手动消息发送属于聊天窗的独立操作。人工内测页则面向已配置样本和人工评价，其结果与正式 Runtime 预览分别记录。

配置以账号和会话为作用域，模型档位、推理强度、回答长度、回复强度、上下文预算和群聊风格分别设置。页面区分管理员偏好、允许范围与锁定状态，并展示本轮有效选择；“上下文长度”指历史读取预算，不能解释成模型服务商的最大上下文窗口。最终使用值以 Runtime 返回结果为准。

\subsection{模型、工具与服务管理}
模型管理按 Luna、Terra、Sol 三个档位组织连接和 Key 池，显示真实模型 ID、接入地址、推理参数、超时、输出额度及 Key 可用状态。相同底层模型可以服务不同档位，选项使用档位标识区分，避免按模型名称合并后选择错位。

配置编辑、保存和运行生效有各自的反馈。模型连接保存后，页面读取 Runtime 配置状态；管理员执行“应用到 Runtime”后，等待服务端确认，再用一次请求观察实际模型。这样可以定位“内容已保存，但运行侧仍使用旧配置”的问题。

\begin{table}[H]\centering
\caption{配置对象与生效反馈}
\begin{tabularx}{\textwidth}{@{}L{28mm}Y Y@{}}\toprule
配置对象 & 页面操作 & 需要观察的结果\\\midrule
群行为与插件 & 选定会话，设置参与方式、偏好和插件范围并保存 & 当前群配置、可用能力、本轮实际选择及启用原因。\\
模型与 Key 池 & 编辑档位连接，维护 Key，检查后应用到 Runtime & 待应用或已应用状态、运行侧是否就绪、回答实际使用的模型。\\
MCP 与扩展插件 & 查看目录，检查连接，按部署能力安装或配置 & 服务连接结果、工具目录与错误说明；是否进入 Planner 另由能力映射决定。\\
群服务档案 & 在控制面选择档案，预览并提交配置 & 当前服务分配、命令结果与后端可用状态。\\\bottomrule
\end{tabularx}
\end{table}

MCP 管理界面提供服务检查与接入，插件管理提供安装和运行状态入口。服务出现在目录中表示已经登记，能否参与智能体回答还取决于当前能力映射、授权与运行配置。群服务页和历史内测页若缺少后端或样本，会展示配置说明及重新检测入口。

\subsection{连接状态、异常与可用性}
WebUI 分别呈现网关连接、QQ 账号在线状态、Agent 预览可用性和生产自动参与状态。网关可访问时，QQ 仍可能需要重新登录；预览接口可用时，自动参与也可能尚未开启。这些状态决定管理员下一步应检查账号、模型还是群策略。

\begin{table}[H]\centering
\caption{试用问题如何转化为交互设计}
\begin{tabularx}{\textwidth}{@{}L{34mm}Y@{}}\toprule
实际问题 & 对应交互与实现处理\\\midrule
首屏等待联系人刷新 & 先显示工作区快照，后台刷新目录；失败时保留已读取内容并提供重试。\\
QQ 离线仍显示连接正常 & 显示在线账号数与账号总数，离线时给出登录指引和重新检测。\\
Agent 预览失败原因不明 & 区分历史读取与 Runtime 调用失败；在对话及运行页展示原因和重试。\\
目录失败被误认为空列表 & 分开表示未读取、读取失败与真实空结果；旧数据标记为上次读取结果。\\
搜索范围和焦点不清楚 & 标注浏览器缓存范围，区分未搜索、加载和无匹配；完善键盘与焦点恢复。\\\bottomrule
\end{tabularx}
\end{table}

默认工作台入口绑定本机回环地址；对外部署由反向代理提供登录保护，群服务管理另外依赖操作员会话与后端授权。页面仅展示管理员完成当前操作所需的信息，模型 Key、OneBot 凭据和私有文件路径由服务端处理。

\subsection{典型操作流程与实现依据}
以启用选课辅助为例：管理员选择 QQ 账号和目标群，确认账号及 Runtime 状态，开启已授权的评课服务与合适的参与方式；随后在 Agent 对话中输入代表问题，在运行页检查实际模型、评课能力、上下文与资料来源；若回答篇幅或参与节奏不合适，再返回配置页调整。工作区把“看到讨论—设置服务—检查回答”组织为连贯操作，支持校园场景中的持续维护。

本章以当前页面、路由、网关与 Runtime 适配器源码为主要依据，并结合开放式试用修复记录。旧版 Web README 中有关早期候选模式和服务数量的描述仅作为历史材料；本章不以其替代当前调用路径。Web 的已有回归结果见第\ref{sec:evaluation}节。
\source{[S16, S18]；前端组件、路由、缓存、网关与 Runtime 适配器，以及开放式试用修复记录。完整路径见附录。}
